\chapter{The Telegram bot}
\label{ch:bot}

\section{Why a bot}

The web application works on a phone, but opening a browser, logging in and
filling a form is still friction, and friction is what kills the habit of
recording expenses. A chat is already open on the phone, so the bot removes the
last excuse: the user types the same lines they used to write in a notes app,
and they land in the database.

\section{No extra dependency}

The bot talks to the Telegram Bot API over plain HTTP with \code{httpx}, which
the project already uses. There is no bot framework.

\begin{lstlisting}[style=py]
API = "https://api.telegram.org/bot{token}/{method}"

async def send_message(chat_id, text, reply_markup=None) -> None:
    payload = {"chat_id": chat_id, "text": text, "parse_mode": "HTML",
               "disable_web_page_preview": True}
    if reply_markup is not None:
        payload["reply_markup"] = reply_markup
    async with httpx.AsyncClient(timeout=TIMEOUT) as c:
        await c.post(_url("sendMessage"), json=payload)
\end{lstlisting}

The API is small: send a message, edit a message, answer a callback, send a
document. Four functions cover everything the application does, and adding a
framework to wrap four HTTP calls would have added a dependency, a second event
loop model and a second way to configure things.

\section{Webhook instead of polling}

There are two ways to receive updates. Polling asks Telegram for new messages in
a loop, which needs a process that is always awake. A webhook lets Telegram POST
each update to a URL.

The application uses a webhook, because on a free host nothing is guaranteed to
stay awake, and an incoming request is exactly what wakes it up.

\begin{lstlisting}[style=py]
async def setup_webhook() -> None:
    if not settings.public_base_url:
        return
    url = settings.public_base_url.rstrip("/") + "/tg/webhook"
    payload = {"url": url, "drop_pending_updates": True,
               "allowed_updates": ["message", "callback_query"]}
    if settings.telegram_webhook_secret:
        payload["secret_token"] = settings.telegram_webhook_secret
    ...
\end{lstlisting}

\code{drop\_pending\_updates} avoids replaying a backlog after a redeploy, and
\code{allowed\_updates} keeps Telegram from sending update types the application
ignores.

\section{Two locks on the door}

The webhook URL is public, so anybody can POST to it. Two independent checks
protect it.

\textbf{First, the secret token.} It is registered with the webhook, and
Telegram sends it back in a header on every update. The router compares it in
constant time:

\begin{lstlisting}[style=py]
secret = settings.telegram_webhook_secret
if secret and not secrets.compare_digest(x_telegram_bot_api_secret_token or "", secret):
    return Response(status_code=403)
\end{lstlisting}

\textbf{Second, the chat allowlist.} Even a genuine Telegram update is only
served if it comes from an allowed chat, because anybody who finds the bot name
can message it.

\begin{lstlisting}[style=py]
def is_allowed(chat_id) -> bool:
    allow = settings.allowed_chat_ids
    return (not allow) or (str(chat_id) in allow)
\end{lstlisting}

An empty allowlist means setup mode: the bot replies with the chat id so the
owner can copy it into the configuration. Once the allowlist has a value,
anything else is ignored in silence, with a line in the log. Answering an
intruder only confirms that the bot exists.

\begin{pitfall}[Never answer 5xx to Telegram]
The webhook returns \code{200} even when handling the update failed. Telegram
retries on server errors, and a bug that raises on one message would turn into
an endless retry of the same message. The error is logged; the update is
acknowledged.
\end{pitfall}

\section{Confirm before writing}

Free text is read by the same parser the web import uses. The result is not
saved: it is shown as a preview with two buttons.

\begin{lstlisting}[style=py]
token = secrets.token_urlsafe(6)
_PENDING[token] = (chat_id, parsed)
while len(_PENDING) > _MAX_PENDING:
    _PENDING.popitem(last=False)          # drop the oldest
...
await send_message(chat_id, body, reply_markup=_kb([
    [("Confirm", f"c:ok:{token}"), ("Cancel", f"c:no:{token}")]]))
\end{lstlisting}

Pressing a button sends a callback with that token, and only then are the rows
inserted. The message is edited in place with the result, so the chat shows what
happened instead of growing a new message.

Pending previews live in a capped \code{OrderedDict} in memory. That is a
deliberate choice with two consequences: a restart loses them, and an old
preview is evicted after fifty newer ones. Both are safe, because losing a
pending preview means nothing was written. Storing them in the database would
have meant a table, a cleanup job and a way to expire rows, to protect data that
does not exist yet.

\section{Periods and ranges}

Several commands accept a period. The parser handles a month, a year, or a range
of either, and returns a label for the report title.

\begin{lstlisting}[style=py]
def _parse_range(arg: str) -> tuple[int, int, int, int, str]:
    toks = [t for t in (arg or "").split() if t.lower() not in _CONNECTORS]
    if not toks:
        y, m = _parse_period("")
        return y, m, y, m, f"{MONTHS[m]} {y}"
    if len(toks) == 1:
        y, m = _parse_period(toks[0])
        if m:
            return y, m, y, m, f"{MONTHS[m]} {y}"
        return y, 1, y, 12, str(y)          # a year alone means the whole year
    ya, ma = _parse_period(toks[0])
    yb, mb = _parse_period(toks[-1])
    return ya, (ma or 1), yb, (mb or 12), f"{MONTHS[ma or 1][:3]} {ya} – ..."
\end{lstlisting}

Connector words such as \emph{to} and \emph{until} are dropped, so
\code{/report 01/2026 to 03/2026} and \code{/report 01/2026 03/2026} both work.
A year without a month opens at January and closes at December, which is what a
person means when they type a year.

The report command then picks the right generator: a single month gives a
monthly report, a full year gives a yearly one, anything else gives a range
report.

\section{Commands}

\begin{table}[h]
\centering
\begin{tabular}{@{}p{4.6cm}p{9.1cm}@{}}
\toprule
\textbf{Command} & \textbf{Answer} \\
\midrule
\code{/menu} & Button menu, which is faster than typing on a phone. \\
\code{/summary [mm/yyyy]} & Income, expenses, savings and net worth. \\
\code{/net [period]} & Income, expenses and the net figure, any range. \\
\code{/networth} & Investments and accounts, with profit and return. \\
\code{/banks}, \code{/bank <name> <balance>} & Read or update balances. \\
\code{/expenses [mm/yyyy]} & Expenses by category. \\
\code{/cashflow [yyyy]} & Income against expenses, month by month. \\
\code{/savings [mm/yyyy]} & Savings of the period. \\
\code{/report [period]} & PDF with charts. \\
\code{/detailed [period]} & Full PDF: net worth, categories, cashflow. \\
\bottomrule
\end{tabular}
\caption{The bot command set.}
\end{table}

The \code{/bank} command is the only one that writes without confirmation, since
it overwrites a single number the user just typed and the reply shows the new
total.

\section{Synchronous and asynchronous sending}

The module has two ways to send. The async functions serve the webhook. A pair
of synchronous ones serve the scheduler, which runs in a background thread with
no event loop.

They differ in how they treat failure, and the difference is intentional:

\begin{lstlisting}[style=py]
def send_document_sync(chat_id, data, filename, caption="", mime="application/pdf"):
    r = httpx.post(_url("sendDocument"), ...)
    r.raise_for_status()          # a failed backup must be loud

def send_message_sync(chat_id, text) -> None:
    try:
        httpx.post(_url("sendMessage"), ...)
    except Exception as exc:
        log.warning("send_message_sync failed: %s", exc)   # a lost reminder is fine
\end{lstlisting}

A month-end reminder that does not arrive costs nothing. A backup that fails
silently is worse than having no backup at all, because it produces false
confidence. So one swallows the error and the other propagates it up to the task
endpoint, which turns it into a \code{500} that the cron shows in red.
